\subsection*{MT-M18. Unconditional arithmetic block synthesis and an absolute adapted basis}

Fix a primitive unitary Hecke character \(\eta\) over a number field.
Use the entire completion \(\Xi=\Xi_\eta=P_\eta\Lambda(s,\eta)\)
and its inverse Mellin transform \(\psi=\psi_\eta\) from MT-M12,
but make no GRH or simplicity assumption.  The polynomial \(P_\eta\)
removes the possible pure-norm poles exactly as there.  Let \(Z\)
be the set of distinct zeros of \(\Xi\), with their actual
multiplicities \(m_\rho\).  These are the nontrivial zeros and
lie in \(0\le\operatorname{Re}s\le1\).
Put
\[
 \begin{gathered}
 \mathcal E=\{f\in C^\infty(\mathbb R,\mathbb C):
   p_{N,j}(f)=\sup_x e^{N|x|}|f^{(j)}(x)|<\infty
                    \text{ for all integers }N,j\ge0\},\\
 Q=\mathcal E/\overline{\psi*\mathcal E}.
 \end{gathered}
\]
All constants below concern this one fixed Hecke function; no
uniformity in its field, conductor, or infinity type is asserted.

There are real numbers
\[
 n+\tfrac14<h_n<n+\tfrac34\quad(n\in\mathbb Z)
\]
not equal to any zero ordinate, with the following properties.
Let \(\Omega_n\) be the interior of the rectangle
\([-1,2]+i[h_n,h_{n+1}]\), and orient its boundary
\(\Gamma_n\) positively.  Let \(A(\Omega_n)\) be the Banach
algebra of functions holomorphic in \(\Omega_n\) and continuous
on its closure, with the supremum norm.  Write
\[
 \begin{split}
 I_n&=\{F\in A(\Omega_n):
        F^{(j)}(\rho)=0\text{ for }\rho\in Z\cap\Omega_n,\
                                  0\le j<m_\rho\},\\
 J_n&=A(\Omega_n)/I_n,\qquad
 \|a\|_n=\inf_{F\bmod I_n=a}\sup_{\Gamma_n}|F|,\qquad
 d_n=\dim_{\mathbb C}J_n=\sum_{\rho\in Z\cap\Omega_n}m_\rho .
 \end{split}                                                    \tag{M18.1}
\]
An empty block has \(J_n=0\).  Then
\[
 d_n=O(\log(2+|n|)),\qquad
 s(J)=\{(a_n):q_k(a):=\sup_n(2+|n|)^k\|a_n\|_n<\infty
                                  \text{ for every }k\ge0\}      \tag{M18.2}
\]
is a Fr\'echet space, and the Mellin block restriction gives a
topological isomorphism
\[
 Q\simeq s(J).                                                    \tag{M18.3}
\]
In particular
\[
 \overline{\psi*\mathcal E}
 =\{f\in\mathcal E:(\mathcal Mf)^{(j)}(\rho)=0
                         \text{ for every }\rho,\ 0\le j<m_\rho\}.
                                                                  \tag{M18.4}
\]
The block expansion in \(Q\) converges absolutely in every
continuous seminorm.  Choosing an Auerbach basis in each
\((J_n,\|\cdot\|_n)\) gives an absolute Schauder basis of \(Q\)
with block-adapted continuous coefficient functionals.

Neither (M18.3) nor this last basis assertion prescribes the
individual zero jets as its coordinate functionals.  Thus the
result does not assert the quantitative condition of MT-M16.
The block partition and its adapted bases are choices; the
closed quotient in (M18.3) is the canonical one.

\begin{proof}
\emph{Analytic inputs and the transform topology.}
NT-Z10 proves unconditionally, including for a fixed
infinite-order unitary character, the polynomial upper bound
for the finite \(L(s,\eta)\) on every fixed vertical strip and
the bound \(O(\log(2+|T|))\) for the number of nontrivial
zeros, with multiplicities, in any fixed-height window.
It also gives
\[
                   |L(2+iT,\eta)|\ge\zeta_K(2)^{-1}.              \tag{M18.5}
\]
The standard completed Hecke functional equation locates all
the zeros of the entire completion in the fixed critical strip.
Exact classical source locators are Ramakrishnan--Valenza,
\emph{Fourier Analysis on Number Fields}, \S7.4,
Theorem 7-19, printed p.~279, and \S7.7, Theorem 7-28,
printed p.~289; the fixed-strip and local-count deductions
are proved in NT-Z10.  We use these unconditional facts,
not a zero-location consequence of either GRH input.

Put \(\kappa=\pi[K:\mathbb Q]/4\).  Stirling's formula and
the preceding strip bound give, for every fixed \(A>0\),
\[
 |\Xi(\sigma+it)|\le C_A(2+|t|)^{B_A}e^{-\kappa|t|}
                         \quad(|\sigma|\le A).                  \tag{M18.6}
\]
Cauchy's formula gives the same type of bound for \(\Xi'\)
on a smaller strip.  For large \(|t|\), the factor
\(\Xi(s)/L(s,\eta)\) is nonzero and has upper and lower
bounds by \(e^{-\kappa|t|}\) times fixed positive or negative
powers of \(2+|t|\), uniformly on each fixed strip.  Its
finitely located gamma exceptions are immaterial at large
height.

For clarity, the transform fact used below is elementary.
If an entire function \(H\) decreases faster than every
power of \(|t|\), uniformly on each strip \(|\sigma|\le A\),
its inverse Fourier integral on \(\sigma=0\) is the Mellin
transform inverse of an element of \(\mathcal E\).
For \(x\ge0\), shift the inverse contour to
\(\sigma=N+1\), and for \(x\le0\) shift it to
\(\sigma=-N-1\).  The horizontal edges tend to zero by the
uniform strip decay.  The \(j\)th derivative is bounded
by the integral of \(|s|^j|H(s)|\) on the shifted line,
times \(e^{-(N+1)|x|}\); these integrals are finite.
This proves all \(p_{N,j}\) bounds and continuity from the
corresponding strip seminorms.  Fourier inversion proves
the transform identity.  Conversely, integration by parts
in \(\mathcal Mf(\sigma+it)=\int e^{(\sigma+it)x}f(x)\,dx\)
shows rapid vertical decay, with each strip seminorm
bounded by finitely many \(p_{N,j}\).  Entireness follows
by exponential domination.  In particular (M18.6) proves
\(\psi\in\mathcal E\) and \(\mathcal M\psi=\Xi\), without GRH.

\emph{Selecting all horizontal boundaries at once.}
For \(n\in\mathbb Z\), let \(M_n\) count, with multiplicities,
the zeros whose ordinates lie in \([n-2,n+3]\).  Thus
\(M_n\le C\log(2+|n|)\).  On
\(I_n^*=(n+1/4,n+3/4)\), set
\[
 U_n(h)=\sum_{\substack{\rho\in Z\\
                    \operatorname{Im}\rho\in[n-2,n+3]}}
       m_\rho\log^+\frac1{|h-\operatorname{Im}\rho|}.
\]
Every summand has integral at most
\(\int_{-1}^1-\log|u|\,du=2\) over \(I_n^*\).
Averaging over its length \(1/2\) therefore permits a
choice \(h_n\in I_n^*\), avoiding every zero ordinate,
with \(U_n(h_n)\le 5M_n\); if \(M_n=0\), use its midpoint.
(The harmless slack from \(4M_n\) to \(5M_n\) gives a
positive-measure set when \(M_n>0\).)
Every zero with \(|h_n-\operatorname{Im}\rho|<1\)
is included in this sum.  Consequently
\[
 \prod_{\rho\in Z}
   \min\{1,|h_n-\operatorname{Im}\rho|\}^{m_\rho}
      \ge e^{-5M_n}\ge c(2+|n|)^{-C}.                            \tag{M18.7}
\]
Only finitely many factors in this product differ from one.
Also \(1/2<h_{n+1}-h_n<3/2\), so the rectangles have
uniformly bounded perimeter and cover all zero ordinates.

We claim, for every fixed \(A>0\), that
\[
 \min_{|\sigma|\le A}|\Xi(\sigma+i h_n)|
       \ge c_A(2+|n|)^{-C_A}e^{-\kappa|n|}.                      \tag{M18.8}
\]
Here is a complete minimum-modulus argument.  For large
\(|h_n|\), center a disk at \(c_n=2+i h_n\) and take a
zero-free boundary of radius
\(R\in[4(A+3)+4,4(A+3)+5]\).  Such a radius exists because
there are only finitely many zeros in that compact annulus.
At sufficiently large height the disk contains no finite-\(L\)
pole or trivial zero, since their possible ordinates are fixed
by the infinity type.  Its total zero multiplicity is
\(M=O_A(\log(2+|n|))\), by NT-Z10.  Remove its zeros by the
finite disk Blaschke product \(B\), with multiplicities,
and write \(L=GB\).  The factor \(G\) is holomorphic and
zero-free on the disk and its boundary.  Since \(|B|=1\)
on the boundary, (M18.6)'s finite-\(L\) antecedent and the
maximum principle give \(|G|\le C(2+|n|)^C\) there and
inside.  At the center, \(|B(c_n)|\le1\), so (M18.5)
gives \(|G(c_n)|\ge c>0\).
Apply Harnack's inequality to the nonnegative harmonic
function \(\log(C(2+|n|)^C)-\log|G|\).  On the inner
half disk its value is at most three times its value
at the center.  Thus
\[
                        |G(s)|\ge c(2+|n|)^{-C}
\]
on that inner disk.  This form of Harnack follows by
bounding the ratios of the Poisson kernels on concentric
circles; the ratio is \((R+r)/(R-r)\le3\) for \(r\le R/2\).

For \(s=\sigma+i h_n\), \(|\sigma|\le A\), the normalized
coordinate \(z=(s-c_n)/R\) has \(|z|<1/4\).
For a disk zero \(\rho\), with \(a_\rho=(\rho-c_n)/R\),
\[
 \left|\frac{z-a_\rho}{1-\overline{a_\rho}z}\right|
    \ge\frac{|s-\rho|}{2R}
    \ge\frac{\min\{1,|h_n-\operatorname{Im}\rho|\}}{2R}.
\]
Products of omitted factors in (M18.7) are at most one.
It follows that
\[
 |B(s)|\ge(2R)^{-M}
   \prod_{\rho\in Z}
      \min\{1,|h_n-\operatorname{Im}\rho|\}^{m_\rho}
             \ge c_A(2+|n|)^{-C_A}.
\]
Multiplication by \(G\) and by the lower Stirling bound
for \(\Xi/L\) proves (M18.8) at large height.
There are only finitely many remaining \(n\); each
horizontal segment is compact and zero-free for \(\Xi\),
so reducing \(c_A>0\) covers them.

The same disk proof, now centered at \(2+it\), gives a
polynomial lower bound for \(L(\sigma+it,\eta)\) on either
fixed vertical line \(\sigma=-A\) or \(\sigma=A\), for
every \(A\ge2\) and large \(|t|\).  Indeed their distances
from all disk zeros in the critical strip are at least one,
so every normalized Blaschke factor is at least \(1/(2R)\).
There is no logarithmic-distance product to control.
The lower Stirling bound and the compact low-height range
give the corresponding lower bound for \(\Xi\).
The same argument applies to the line \(\sigma=-1\)
and to \(\sigma=2\).  Combining it with (M18.8) yields
\[
 \inf_{w\in\Gamma_n}|\Xi(w)|
       \ge c(2+|n|)^{-C}e^{-\kappa|n|}.                          \tag{M18.9}
\]
All constants here are independent of separations between
zeros inside a rectangle.

\emph{The finite block algebras and an entire interpolation kernel.}
All zeros of a block are interior points.  Each of its
finitely many jet functionals is continuous on
\(A(\Omega_n)\), by Cauchy's formula on a disk contained
in \(\Omega_n\).  The map to these jets is onto, by
ordinary polynomial Hermite interpolation.  Its kernel
is \(I_n\), a closed ideal, proving that (M18.1) is a
genuine finite-dimensional quotient norm, with dimension
\(d_n\).  The maximum principle identifies the supremum
over the closure with that over \(\Gamma_n\).
The quotient norm is submultiplicative because the
supremum norm is.  The local zero count gives (M18.2)'s
dimension bound.

Fix \(a>0\).  For \(F\in A(\Omega_n)\) define
\[
 H_nF(z)=\frac1{2\pi i}\int_{\Gamma_n}
    F(w)\frac{\Xi(z)-\Xi(w)}{\Xi(w)(z-w)}
                      e^{a(z-w)^2}\,dw .                       \tag{M18.10}
\]
The apparent singularity at \(z=w\) is removable.
Consequently this is an entire function of \(z\).
If \(z\) is inside the rectangle, Cauchy's formula
for the second term gives
\[
 H_nF(z)=F(z)+\Xi(z)\frac1{2\pi i}
             \int_{\Gamma_n}
               \frac{F(w)e^{a(z-w)^2}}{\Xi(w)(z-w)}\,dw .
                                                                  \tag{M18.11}
\]
The integral on the right is holomorphic in the interior.
If \(z\) is outside the closed rectangle, that second
Cauchy integral is zero and \(H_nF(z)\) is divisible
locally by \(\Xi(z)\).
These statements follow also for continuous boundary
values by applying the formulas on interior rectangles
and taking limits.  Thus \(H_nF\) matches every
prescribed block jet of \(F\), and vanishes to the
full multiplicity at every zero outside the block.

Moreover, if all the block jets of \(F\) vanish, then
\(F/\Xi\) is holomorphic in \(\Omega_n\) and continuous
on its closure.  For exterior \(z\), both terms in
(M18.10) then integrate to zero by Cauchy's theorem.
Entireness implies \(H_nF=0\) everywhere.  It follows
that \(H_n\) depends only on \(a_n=F\bmod I_n\).
Denote this entire function by \(H_n a_n\).

For every fixed \(A>0\) there are constants \(C_A,B_A'>0\)
and \(c_a>0\) such that
\[
 |H_n a_n(\sigma+it)|
       \le C_A(2+|n|)^{B_A'}\|a_n\|_n
                         e^{-c_a(t-n)^2},
                 \qquad |\sigma|\le A.                         \tag{M18.12}
\]
To check the estimate before taking the quotient infimum,
bound \(|F|\) by its boundary supremum in (M18.10).
If \(|z-w|\ge1\), use (M18.6) and (M18.9) to bound
the difference quotient by
\[
 C_A(2+|n|)^C(2+|t|)^C
                    e^{\kappa\,\lvert t-\operatorname{Im}w\rvert}
       +1 .
\]
Here \(|\operatorname{Im}w-n|<2\), while
\[
 |e^{a(z-w)^2}|
       \le C_A e^{-a(t-\operatorname{Im}w)^2}.
\]
A fixed polynomial and the fixed linear exponential
are absorbed by a smaller Gaussian, and
\(2+|t|\le(2+|n|)(2+|t-n|)\).
If \(|z-w|<1\), the identity
\[
 \frac{\Xi(z)-\Xi(w)}{z-w}
      =\int_0^1\Xi'(w+u(z-w))\,du
\]
and the strip derivative bound give a polynomial in
\(2+|n|\) after division by (M18.9).
In this region \(|t-n|<3\), so it too satisfies (M18.12)
after a constant adjustment.  The perimeter is bounded
independently of \(n\).  Taking the infimum over all
representatives of \(a_n\) proves (M18.12).

The transform fact above now gives a unique
\(\mathcal T_n(a_n)\in\mathcal E\) whose Mellin transform is
\(H_n a_n\).  It is linear in \(a_n\), and for every
defining seminorm
\[
 p_{N,j}(\mathcal T_n(a_n))
       \le C_{N,j}(2+|n|)^{B_{N,j}}\|a_n\|_n .                  \tag{M18.13}
\]
Indeed on either shifted vertical line, the integral
of \((1+|t|)^j\) against (M18.12)'s Gaussian costs
only another fixed power of \(2+|n|\).
In particular no individual reciprocal derivative
or Hermite interpolation matrix norm has entered
these estimates.

\emph{Spectral synthesis for the closed ideal.}
Write \(J\) for the right side of (M18.4).
It is a closed subspace, since each Mellin jet is
continuous.  Multiplication of Mellin transforms
immediately gives
\(\overline{\psi*\mathcal E}\subseteq J\).
For the reverse inclusion take \(f\in J\) and
\(F=\mathcal Mf\).  The quotient \(F/\Xi\) is entire:
every zero, including its full multiplicity, has
been cancelled.

Fix \(A\ge2\).  Use the wider rectangles
\([-A,A]+i[h_n,h_{n+1}]\).  On their horizontal
edges (M18.8) applies; on their vertical edges
the vertical-line argument following (M18.8)
applies.  The rapid fixed-strip decay of \(F\)
and these lower bounds show, on the whole boundary,
\[
 \left|\frac{F(s)}{\Xi(s)}\right|
       \le C_{A,f}(2+|n|)^{C_A}e^{\kappa|n|}.
\]
The maximum principle for the entire quotient
extends this bound through the rectangle.
Since their heights have gaps between \(1/2\)
and \(3/2\), it follows that on every fixed strip
\[
 \left|\frac{F(\sigma+it)}{\Xi(\sigma+it)}\right|
       \le C_{A,f}(2+|t|)^{C_A}e^{\kappa|t|}.                  \tag{M18.14}
\]
For every \(\epsilon>0\), the entire function
\[
                   e^{\epsilon s^2}F(s)/\Xi(s)
\]
therefore has rapid decay on every fixed strip.
It is the Mellin transform of some
\(j_\epsilon\in\mathcal E\).
With \(g_\epsilon(x)=(4\pi\epsilon)^{-1/2}
e^{-x^2/(4\epsilon)}\), Fourier uniqueness gives
\[
                         \psi*j_\epsilon=g_\epsilon*f.
\]
The right side tends to \(f\) in \(\mathcal E\).
For completeness, on \(|y|\le1\) the weighted
mean-value bound for
\(f^{(j)}(x-y)-f^{(j)}(x)\) is
\(|y|e^{N|y|}p_{N,j+1}(f)\).
On the complement of any fixed neighborhood of
zero use
\((e^{N|y|}+1)p_{N,j}(f)\).
Integration against \(g_\epsilon(y)\) makes both
bounds tend to zero, by its first moment and
Gaussian tails.  This proves the reverse inclusion
and hence (M18.4), independently of the interpolation
construction.

\emph{The topological decomposition and the adapted basis.}
The restriction map
\[
 R:\mathcal E\longrightarrow s(J),\qquad
        (Rf)_n=(\mathcal Mf)|_{\Omega_n}\bmod I_n
\]
is continuous.  Indeed the quotient norm is at most
\(\sup_{\Gamma_n}|\mathcal Mf|\), and these boundary
suprema decrease faster than every power of \(|n|\),
with each such bound controlled by finitely many
seminorms of \(f\).  Define
\[
                    S(a)=\sum_{n\in\mathbb Z}\mathcal T_n(a_n).
\]
For each seminorm \(p_{N,j}\), (M18.13) and
\(q_{B_{N,j}+2}(a)\), with the exponent rounded up
if necessary, bound the sum of its term seminorms
by a constant times \(\sum_n(2+|n|)^{-2}\).
Thus this series converges absolutely in every
seminorm of \(\mathcal E\), and defines a continuous
linear map \(S:s(J)\to\mathcal E\).
Completeness of \(s(J)\) follows directly from the
complete finite-dimensional norms in its weighted
supremum seminorms; completeness of \(\mathcal E\)
follows from convergence of all weighted derivative
seminorms on compact intervals.

Every Mellin jet of \(S(a)\) may be computed
termwise by continuity.  The interpolation property
of (M18.11) gives \(RS=\mathrm{id}\).
Also \(\ker R=J=\overline{\psi*\mathcal E}\).
Therefore \(R\) induces the continuous bijection
(M18.3), with continuous inverse
\(a\mapsto[S(a)]\).
The block operators \(\mathcal T_n R_n\) on \(\mathcal E\)
are finite-rank continuous projections, mutually
annihilating for distinct \(n\).
Their absolute sum is the continuous projection
\(SR\); on \(Q\) their sum is the identity.

Finally every finite-dimensional complex normed
space has an Auerbach basis \(v_{n,1},\ldots,v_{n,d_n}\):
its vectors and dual coordinate functionals
\(\phi_{n,j}\) all have norm one.
Here is the elementary existence argument.  In
any fixed coordinates maximize the absolute
determinant of \(d_n\) vectors in the compact unit
ball.  A maximum is positive.  Replacing its
\(j\)th column by any unit-ball vector multiplies
the determinant by the corresponding dual
coordinate.  Maximality makes that coordinate
have modulus at most one; its value on its own
column is one.  This proves the claimed norms.

For \(a=(a_n)\in s(J)\), its resulting scalar
coefficients satisfy
\(|\phi_{n,j}(a_n)|\le\|a_n\|_n\).
Writing \(e_{n,j}\) for \(v_{n,j}\) in its block
and zero elsewhere, one has
\[
 \sum_{n,j}q_k\bigl(\phi_{n,j}(a_n)e_{n,j}\bigr)
 \le\sum_n d_n(2+|n|)^k\|a_n\|_n
 \le q_{k+3}(a)\sum_n d_n(2+|n|)^{-3}<\infty .
\]
This proves an absolute expansion in every
seminorm, in any enumeration of these vectors.
The coordinate functionals are continuous and
unique by the block isomorphism and the finite
basis in each block.  Transport through (M18.3)
proves the asserted absolute Schauder basis.
\end{proof}

The spectral block \(J_n\), as an algebra, is the
finite-length complex algebra
\[
 \mathbb C[s]\Big/
       \prod_{\rho\in Z\cap\Omega_n}(s-\rho)^{m_\rho}
 \simeq
 \prod_{\rho\in Z\cap\Omega_n}
                  \mathbb C[z]/(z^{m_\rho}).
\]
Its stable norm is the quotient holomorphic supremum
norm in (M18.1), not a norm obtained by declaring its
individual jet coordinates orthonormal.
Thus the construction is a finite-length spectral
divisor decomposition compatible with the Mellin
convolution algebra.  It directly refines the repair
of Masullo's Lemma 9.3.2.5, original PDF pp.~425--426:
closure of the convolution ideal is essential, and
the adapted basis has different coordinates from
that lemma's prescribed jets.  The inherited programme
locators are cote~1's nuclear-trace framework and
SA23 geo-H-lidskii/geo-H-genus, as carried in the
attached returns; no new inspection of original
cote scans is claimed here.  Nuclearity and the
trace action of these finite spectral algebras
are treated separately in MT-M19.

No first-priority claim is made for spectral synthesis in related
function spaces; the explicit block norms and bounds above supply
the precise topology and interpolation statement used here.

